% 19 November 2026. Prepared for Edward Baker with OpenAI Codex assistance.
% Lead recorded gpt-6-astra/max. Program statements are prospective.
\section{Synthesis, controls and outstanding challenges}
\label{sec:program}

\subsection{The analytical information being carried forward}

The pair coordinates compress a local divisor without discarding its
complete exterior. The gap chain makes this exterior geometric: nearby
gaps strongly influence the response, while distant zeros enter through
summable dipole tails. The integer hierarchy keeps the shared divisor
constraints; the complex Mellin extension supplies another representation
with its own convergence and branch requirements. These are compatible
views of one physical field. An arbitrary vector of moments need not
come from that field, and evolving a truncated head without its boundary
rows need not preserve that compatibility.

The earlier manuscript proves a common-time approximation by initializing
one finite paired product and evolving it with the polynomial heat
operator \cite[\S\S6--8]{VCDCMManuscript}. Its fixed-input convergence,
counted-contour transfer and conditioning qualifications are relevant
background. The present program uses the arithmetic approximation as a
more structured candidate reference, rather than assuming that generic
product convergence already supplies a useful relative error at large
height. Both require a count, a denominator floor and a complete path.

The normalization issue has two different manifestations. If
$H=AQ$ locally with $A$ analytic and nonzero, a contour integral of
$\partial_z\log A$ contributes no divisor count. But
\[
 \frac{H_{zz}}{H_z}(a)
   =\frac{Q_{zz}}{Q_z}(a)+2\frac{A_z}{A}(a)
\]
at a simple zero $a$ of $Q$. Thus the same factor affects zero velocity
under the heat law. A spatially constant multiplier cancels from that
ratio; an evolving reference with such a multiplier still has to pay
its time defect. These observations explain the earlier localization
analogy without constructing a localization or index-theoretic proof.
A count of two is the same before and at a double zero, so count
invariance alone cannot protect a gap
\cite{VCDCMDiscussion07,VCDCMDiscussion08}.

\subsection{A common contour makes the comparison explicit}

A short derivation specifies how a field estimate becomes a gap estimate.
Let a positively oriented contour $\Gamma$ contain exactly two zeros,
with multiplicity, of each of two holomorphic fields $H$ and $P$, and
no zeros on $\Gamma$. Write their centers and discriminants as
$c_H,D_H$ and $c_P,D_P$. Suppose $H/P$ has a single-valued analytic
logarithm $h$ on a neighborhood of the contour, for example because
$|H-P|<|P|$ throughout that neighborhood. For the power sums
\[
 A_j(H)=\frac1{2\pi\ii}\oint_\Gamma z^j\frac{H_z}{H}\dd z,
 \quad c_H=A_1(H)/2,\quad D_H=A_2(H)/2-c_H^2,
\]
integration by parts gives
\begin{equation}
 A_j(H)-A_j(P)=-\frac{j}{2\pi\ii}
                 \oint_\Gamma z^{j-1}h(z)\dd z.
 \label{prog:log-moments}
\end{equation}
If
$I_0=(2\pi\ii)^{-1}\oint h\dd z$ and
$I_c=(2\pi\ii)^{-1}\oint(z-c_P)h\dd z$, then
\begin{equation}
 c_H-c_P=-\tfrac12 I_0,\qquad
 D_H-D_P=-I_c-\tfrac14 I_0^2.
 \label{prog:centered-error}
\end{equation}
To verify the second equality, substitute
$A_2(H)-A_2(P)=-2(2\pi\ii)^{-1}\oint zh\dd z$ and expand
$c_H^2=(c_P-I_0/2)^2$. The shared logarithm and full square remain.
Separately maximizing their contributions may give a valid but weaker
outer bound. A constant $h$ gives $I_0=I_c=0$ exactly.

Together with the joint quotient identity of
Section~\ref{sec:residual-transport}, this gives a concrete paired
comparison: direct contour evaluation and time-integrated drift must
use the same source information and maintain the same correlations.
These identities rederive the relevant transfer principle in
\cite[\S\S7,19--20]{VCDCMManuscript}. They are not evidence that an
iteration of the same constraints adds information.

\subsection{What a first research result should accomplish}

The first comparison should fix one reference field, one parameter domain
and one target functional before numerical refinement. For a real pair,
that functional may be the cumulative drift in
Proposition~\ref{cont:cumulative}. An alternative is the whole-time
forcing floor in Theorem~\ref{cont:barrier}. Their acquisition requirements
are different, and an estimate for one must not be silently substituted
for the other.

There are three connected analytical possibilities. First, estimate the
joint residual numerator directly so that common multiplicative errors
cancel before an absolute value is taken. Second, use the forced heat
representation or its adjoint to weight only those residual components
which reach the selected contour functional, while paying the full
spatial complement and boundary data. Third, retain the phases of the
arithmetic correction and its reference sum so that the corrected center
is bounded as one expression. These are proposed mechanisms, not acquired
canonical bounds. A mollifier that improves a value bound is not
automatically a stable reference for a derivative or time-dependent
quantity.

A successful finite result would be a smaller complete error allowance
or a longer certified interval for the fixed consumer, compared with
a direct method using the same information. A failure could also be
useful if it identifies the particular discarded correlation or growing
constant that prevents improvement. Such a failure is scoped to that
comparison. It cannot justify a universal claim that dynamics, different
mollifiers or other arithmetic decompositions are incapable of progress.

\subsection{Collision controls and failed premises}

The collision controls in \cite{RepositoryControls} distinguish a valid
transfer from an unjustified source-specific conclusion. The following
is an analytical premise screen, not a numerical replay of the control
certificates.

The simplest exact test is
\begin{equation}
 H(t,z)=P(t,z)=z^2-2(t-T),\qquad T>0.
 \label{prog:quadratic}
\end{equation}
It has zero approximation error and zero heat defect, but its real pair
merges at $t=T$. Here $M=K=0$, $V=1$ and
$D(u)=2(t_*-T)-2u$. The cumulative allowance is exhausted exactly at the
merger; the strict margin fails. Thus a sound residual theorem preserves
the event. Perfect approximation of a colliding reference is not
protection against collision.

An initial almost regular chain also need not provide a whole-time
forcing floor. For $k>0$ and $0<\alpha_0<1$, the exact backward heat
field
\begin{equation}
 F(u,z)=e^{-k^2u}\cos(kz)+\alpha_0
 \label{prog:cosine-control}
\end{equation}
has $F_u=F_{zz}$. Put $\alpha(u)=\alpha_0e^{k^2u}$ and
$\theta(u)=\arccos\alpha(u)$. Its short pair about $\pi/k$ has
\[
 D(u)=\theta(u)^2/k^2,\qquad
 V(u)=\frac{\theta(u)\alpha(u)}{\sqrt{1-\alpha(u)^2}},\qquad
 u_c=k^{-2}\log(1/\alpha_0).
\]
The formulas follow by solving $\cos(kz)=-\alpha(u)$ and differentiating.
As $\alpha_0\to0$, the initial gaps approach $\pi/k$ and
$V(0)\to0$, yet a merger occurs at $u_c$. For fixed $k$ the exterior
of the selected cluster stays separated there; its inverse powers
remain summable, so $K\to0$ as $D\to0$. The whole-time near-unit
forcing premise fails. This is a finite positive atomic source, not
theta, and serves only as a control on inference from initial geometry.

The remaining source controls fail different canonical premises:
\begin{itemize}
\item Euler-factor translates and weighted-square, finite-source,
finite-score and finite-prime perturbations alter the full normalized
theta arithmetic. Preserved finite measurements do not restore its
complete source identity. Documented endpoint-defective subclasses also
lack the imported strip.
\item Reflected super-Gaussian sources fail the order-one product input
used here as well as literal theta identification. Smooth shape and
positive-primitive controls lack that complete arithmetic identification;
positivity alone supplies no missing gap margin.
\item Atomic and finite-profile constructions fail the full theta profile.
Their event does not refute a theorem that uses the actual profile, but
it must survive any generic approximation theorem whose hypotheses they
satisfy.
\item RH-conditional endpoint constructions change ordinary Dirichlet
arithmetic and assume RH in their construction. Their escaping event
heights do not provide a fixed-compact obstruction.
\item Changed observation kernels fail the physical cosine/Newman
readout; fractional profiles fail integer divisor multiplicities; static
divisor surgery supplies no common source heat path. Strong-collar and
annular polynomial controls still test local algebra and the full square.
\end{itemize}
Identifying a missing canonical premise is not proof that a proposed
coarse estimate uses that premise effectively. Every future quantitative
claim must identify where the unchanged arithmetic actually enters.

\subsection{Outstanding analytical challenges}

The remaining work has a definite dependency order.

\begin{enumerate}
\item \emph{Acquire a useful reference and initial gap.}
The full arithmetic approximation is available beyond first-term clock
dominance, but its small remainder does not locate its own zeros.
A counted contour, initial separation and a compatible correction center
must be established without assuming the desired continuation.

\item \emph{Control the residual in the norm the gap uses.}
The forced error equation needs time-uniform spatial derivatives and
complete tail or boundary data. Its inverse denominator and center
errors must be paid. The main question is whether the common numerator,
its weighted time integral, or arithmetic phase information produces a
useful improvement after all of those costs.

\item \emph{Keep the changing environment coupled.}
The Mellin hierarchy organizes higher moments but is not a finite closure.
A source-constrained realization or direct source estimate must retain
shared zeros, moving gaps and complete exterior forces. At a regular
merger $K=2D^2\mathcal F$ vanishes, so assuming a whole-path positive
floor without acquiring it would already assume event exclusion.

\item \emph{Extend the source domain, not just its precision.}
At fixed positive time the favorable moment estimates inherit individual
clock separation. A new result must state exactly which time/height
domain it adds or which complete comparison it improves. The onset
$x\ge\exp(C/t)$ is a sufficient source threshold; its divergence at zero
is not removed by substituting an unevaluated constant or taking a
nonuniform limit.

\item \emph{Assemble local information into a global statement.}
A protected pair or finite block must be joined to compact and high-height
regions with no omitted gaps, interfaces or exits. If the Polymath strip
criterion is the eventual consumer, its initial, exterior and
intermediate-time barrier hypotheses and its $y_0^2/2$ time payment
remain. A local real-gap estimate cannot replace those obligations
without a new proved bridge.
\end{enumerate}

The purpose of the new track is to resolve these analytical dependencies
in sustained research sequences: derive a candidate estimate, test its
premises, compare its complete consequence, and use the result to choose
the next step. Numerical work supports those decisions through fixed
algebraic checks or explicitly certified examples. The unresolved
mathematical task is a source-specific estimate whose useful margin
persists beyond the information already supplied by direct localization.
